\documentclass[11pt]{article}
\usepackage[margin=1in]{geometry}
\usepackage{amsmath,amssymb}
\usepackage{enumitem}
\usepackage{hyperref}
\usepackage{longtable}
\usepackage{array}
\usepackage{listings}
\usepackage{xcolor}

\lstset{
  basicstyle=\ttfamily\small,
  breaklines=true,
  columns=fullflexible,
  frame=single,
  keepspaces=true,
  keywordstyle=\color{blue!60!black},
  commentstyle=\color{green!40!black},
  stringstyle=\color{red!50!black},
  showstringspaces=false
}

\title{2022 Paper to Code Mapping\\Magnetic-Field-Inspired Navigation}
\author{}
\date{2026-06-18}

\begin{document}
\maketitle

\section*{Purpose}

This document compares the 2022 paper
\emph{Magnetic-Field-Inspired Navigation for Robots in Complex and Unknown Environments}
with the current implementation in
\texttt{src/mfinav/double\_integrator.py}.

The goal is not to claim equivalence. The goal is to make the current status
explicit:

\begin{itemize}[leftmargin=*]
  \item what the paper says,
  \item where that idea appears in the code,
  \item whether the implementation is exact, approximate, or missing.
\end{itemize}

Primary paper source:
\url{https://www.frontiersin.org/journals/robotics-and-ai/articles/10.3389/frobt.2022.834177/full}

\section*{Legend}

\begin{itemize}[leftmargin=*]
  \item \textbf{Exact}: the current code follows the same structure and intent closely enough to be treated as the same mechanism.
  \item \textbf{Approximate}: the current code is clearly inspired by the paper but does not yet implement the same equation or assumption literally.
  \item \textbf{Missing}: the paper concept is not yet represented faithfully in code.
\end{itemize}

\section*{Reference Code Locations}

\begin{itemize}[leftmargin=*]
  \item Goal controller: lines 225--286
  \item Local sensing and closest-point averaging: lines 289--327
  \item Boundary-following field: lines 330--357
  \item Collision-avoidance field: lines 360--376
  \item Combined control law: lines 379--396
  \item Simulation loop and saturation assumptions: lines 567--608
\end{itemize}

\section*{High-Level Control Law}

Paper Section 4 states that the obstacle term is split into a
boundary-following vector field and a collision-avoidance vector field. Paper
Section 5 states that the point robot is steered by control law (38), with a
goal term and the obstacle term.

\begin{equation}
u = F_g + F_o
\end{equation}

\begin{equation}
F_o = F_b + F_a
\end{equation}

\textbf{Current code}

The current code follows exactly that high-level decomposition:

\begin{itemize}[leftmargin=*]
  \item \texttt{goal\_cmd} in lines 392--393
  \item \texttt{boundary\_cmd} in lines 394
  \item \texttt{collision\_cmd} in line 395
  \item sum in line 396
\end{itemize}

\textbf{Status}: \textbf{Exact at architecture level}

\section*{Go-To-Goal Term \(F_g\)}

\subsection*{Paper}

Paper Section 5 says the authors first used PD control from Eq. (39) as the
goal attraction term. The same section then says geometric control was
introduced as an alternative because plain PD does not guarantee constant speed.

The PD form described in the paper is the standard point-robot attractive term:

\begin{equation}
F_g = -K_P (p - p_g) - K_D \dot{p}
\end{equation}

Equivalently, with \(e = p - p_g\),

\begin{equation}
F_g = -K_P e - K_D \dot{p}
\end{equation}

\subsection*{Current code}

The current code supports three goal modes:

\begin{itemize}[leftmargin=*]
  \item \texttt{goal\_mode = "pd"} at lines 269--273 and 274--286
  \item \texttt{goal\_mode = "geometric"} at lines 274--284
  \item \texttt{goal\_mode = "hybrid"} at lines 274--286
\end{itemize}

In PD mode, the implemented term is:

\begin{equation}
u_g = k_{\text{scale}} K_P (p_g - p) - K_D \dot{p}
\end{equation}

which is equivalent to the paper PD form when \(k_{\text{scale}} = 1\).

In geometric mode, the current code uses:

\begin{equation}
u_g = -k K_{P,r} (\lVert v \rVert - v_{\max}) \hat{v}
      + K_{\omega} \lVert v \rVert \hat{n} (\hat{n}^{\top}\hat{g})
\end{equation}

In hybrid mode, the geometric-like term is used far from the goal and the PD
term is used near the goal.

\subsection*{Code excerpt}

\begin{lstlisting}[language=Python]
if self.config.use_legacy_goal_relaxation:
    kp_scale = self._goal_weight(goal_vector, observation.closest_obstacle_vector)
else:
    kp_scale = 1.0

if self.config.goal_mode == "pd":
    return kp_scale * self.config.kp_goal * goal_vector - self.config.kd_goal * state.velocity

if self.config.goal_mode == "geometric":
    speed = _norm(state.velocity)
    if goal_mag > self.config.magni_bound:
        if speed > EPS:
            direction = _unit(state.velocity)
        else:
            direction = _unit(goal_vector)

        speed_error = -(speed - self.config.speed_limit)
        vel_attr = (self.config.kp_goal_relaxed * kp_scale) * speed_error * direction

        if speed > EPS:
            angle_error = _signed_angle(state.velocity, goal_vector)
            omega = self.config.kp_geom * kp_scale * angle_error
            force_add = omega * _perp_left(state.velocity)
        else:
            force_add = np.zeros(2, dtype=float)
        return vel_attr + force_add

    return kp_scale * self.config.kp_goal * goal_vector - self.config.kd_speed * state.velocity

return kp_scale * self.config.kp_goal * goal_vector - self.config.kd_speed * state.velocity
\end{lstlisting}

This is closer to the ROS-era SO(3) geometric steering than the earlier
simplified lateral heuristic, but it is still a 2D analogue rather than a full
matrix-logarithm reconstruction.

\subsection*{Comparison}

\begin{longtable}{|p{0.23\linewidth}|p{0.24\linewidth}|p{0.14\linewidth}|p{0.29\linewidth}|}
\hline
\textbf{Paper item} & \textbf{Code location} & \textbf{Status} & \textbf{Notes} \\
\hline
PD goal term \(F_g\) from Eq. (39) & lines 269--273, 286 & Approximate & Correct sign and structure in PD mode, but optional scaling \(k_{\text{scale}}\) can modify the pure paper form. \\
\hline
Geometric control alternative & lines 274--284 & Approximate & Present in spirit, but not verified line-by-line against the geometric controller referenced by the paper. \\
\hline
Implementation choice ``PD first, geometric alternative'' & lines 200--223, 269--286 & Approximate & The repo exposes explicit modes, but the default pragmatic path still uses a hybrid controller not stated in the paper. \\
\hline
\end{longtable}

\section*{Boundary-Following Field \(F_b\)}

\subsection*{Paper}

Paper Section 4.1 says:

\begin{itemize}[leftmargin=*]
  \item \(F_b\) is induced by an artificial current on the obstacle surface,
  \item the artificial current \(l_o\) has the same direction as the
  projection of the robot motion direction \(l_a\) onto the obstacle surface,
  \item \(F_b\) is the term responsible for boundary following.
\end{itemize}

The paper also states in Lemma 1 that \(F_b\) does not affect robot speed.

\subsection*{Current code}

Current code computes:

\begin{equation}
\hat{l}_a = \frac{v}{\lVert v \rVert}
\end{equation}

\begin{equation}
\hat{r} = \frac{-r_o}{\lVert r_o \rVert}
\end{equation}

\begin{equation}
l_o \approx \hat{l}_a - (\hat{l}_a^{\top}\hat{r})\hat{r}
\end{equation}

with a perpendicular fallback when the projection becomes too small
(lines 342--353).

Two field modes now exist:

\begin{itemize}[leftmargin=*]
  \item \textbf{paper mode} at lines 355--357:
  \begin{equation}
  F_b^{\text{code}} =
  \sigma_b c\, l_{a,\perp}\left(l_{a,\perp}^{\top}
  \left(\frac{\lVert v \rVert}{d} l_{o,\perp}\right)\right)
  \end{equation}
  \item \textbf{pragmatic mode} at lines 356--357:
  \begin{equation}
  F_b^{\text{code}} =
  \sigma_b c\, v_{\text{guide}} \frac{\hat{l}_o}{d}
  \end{equation}
\end{itemize}

\subsection*{Code excerpt}

\begin{lstlisting}[language=Python]
dist_to_obs = observation.used_obstacle_vector
dist_from_obs = -dist_to_obs
dist_from_obs_hat = _unit(dist_from_obs)

obs_cur = agent_cur - float(np.dot(agent_cur, dist_from_obs_hat)) * dist_from_obs_hat
obs_cur_mag = _norm(obs_cur)
if self.config.field_mode == "paper":
    if obs_cur_mag >= self.config.epsilon_current:
        obs_cur = obs_cur / obs_cur_mag
        self.previous_surface_current = obs_cur.copy()
    elif self.previous_surface_current is not None:
        obs_cur = self.previous_surface_current.copy()
    else:
        return np.zeros(2, dtype=float)

    velocity_3 = _embed_2d(state.velocity)
    agent_current_3 = _embed_2d(agent_cur)
    obs_current_3 = _embed_2d(obs_cur)
    b_field_3 = _skew3(obs_current_3) @ velocity_3 / max(distance, EPS)
    force_3 = self.config.c_field * (_skew3(agent_current_3) @ b_field_3)
    return _project_2d(force_3)
\end{lstlisting}

The pragmatic mode still uses the direct tangent-guidance surrogate:

\begin{lstlisting}[language=Python]
if obs_cur_mag < self.config.epsilon_current:
    obs_cur = _perp_left(dist_from_obs_hat)
else:
    obs_cur = obs_cur / obs_cur_mag

sigma_b = max(0.0, min(1.0, (self.config.r_l - distance) / max(self.config.r_l - self.config.r_la, EPS)))
guidance_speed = max(speed, 0.5 * self.config.speed_limit)
return sigma_b * self.config.c_field * guidance_speed * obs_cur / max(distance, EPS)
\end{lstlisting}

\subsection*{Comparison}

\begin{longtable}{|p{0.23\linewidth}|p{0.24\linewidth}|p{0.14\linewidth}|p{0.29\linewidth}|}
\hline
\textbf{Paper item} & \textbf{Code location} & \textbf{Status} & \textbf{Notes} \\
\hline
Artificial current as projection of robot motion onto obstacle surface & lines 342--353 & Exact in spirit & This part is structurally aligned with the paper. \\
\hline
Magnetic boundary-following field derived from that artificial current & lines 355--357 & Approximate & The code uses a 2D surrogate that captures the idea but is not yet verified against the exact paper equation. \\
\hline
Property that \(F_b\) does not affect speed & lines 355--357, 599--600 & Approximate / violated by simulation assumptions & The paper claims this property analytically, but our simulator also applies acceleration and speed clipping. \\
\hline
\end{longtable}

\section*{Collision-Avoidance Field \(F_a\)}

\subsection*{Paper}

Paper Section 4.2 states that \(F_a\) is needed for concave obstacles because
boundary following alone does not regulate the distance to the obstacle
surface. Lemma 5 states that \(F_a\) also does not affect speed.

\subsection*{Current code}

Current code uses:

\begin{equation}
F_a^{\text{code}} =
-\sigma_a c_{\perp}
\left(
\frac{1}{d} - \frac{1}{r_{la}}
\right)
\frac{r_o}{d}
\end{equation}

In \texttt{field\_mode = "paper"}, the code additionally projects out the
component along the motion direction if the robot is moving (lines 372--375).

\subsection*{Code excerpt}

\begin{lstlisting}[language=Python]
if self.config.field_mode == "paper":
    repulsion = -self.config.c_perp * (1.0 / max(distance, EPS) - 1.0 / self.config.r_la) * (
        dist_to_obs / max(distance, EPS)
    )
    if _norm(agent_cur) > EPS:
        repulsion = repulsion - float(np.dot(repulsion, agent_cur)) * agent_cur
    return repulsion

sigma_a = max(0.0, min(1.0, (self.config.r_la - distance) / max(self.config.r_la, EPS)))
repulsion = -sigma_a * self.config.c_perp * (1.0 / max(distance, EPS) - 1.0 / self.config.r_la) * (
    dist_to_obs / max(distance, EPS)
)
if _norm(agent_cur) > EPS:
    repulsion = repulsion - float(np.dot(repulsion, agent_cur)) * agent_cur
return repulsion
\end{lstlisting}

\subsection*{Comparison}

\begin{longtable}{|p{0.23\linewidth}|p{0.24\linewidth}|p{0.14\linewidth}|p{0.29\linewidth}|}
\hline
\textbf{Paper item} & \textbf{Code location} & \textbf{Status} & \textbf{Notes} \\
\hline
Repulsive collision-avoidance field for concave cases & lines 364--376 & Exact in purpose & The code clearly contains the intended mechanism. \\
\hline
Exact paper field expression for \(F_a\) & lines 372--375 & Approximate & The implementation is still a simplified radial repulsion form. \\
\hline
Property that \(F_a\) does not affect speed & lines 372--375 & Approximate & The projection in paper mode tries to preserve the paper property, but this is not yet a verified reconstruction. \\
\hline
\end{longtable}

\section*{Thresholded Use of \(r_l\) and \(r_{la}\)}

\subsection*{Paper}

Paper Section 5 states:

\begin{itemize}[leftmargin=*]
  \item \(F_b\) is only used when the closest distance falls below \(r_l\),
  \item \(F_a\) is only used when the robot gets closer still, below \(r_{la}\).
\end{itemize}

\subsection*{Current code}

\begin{itemize}[leftmargin=*]
  \item \(F_b = 0\) if \(d > r_l\) at lines 335--337
  \item \(F_a = 0\) if \(d > r_{la}\) at lines 365--367
  \item smooth activation weights \(\sigma_b\) and \(\sigma_a\) are used at
  lines 355 and 372
\end{itemize}

\textbf{Status}: \textbf{Approximate}

The threshold logic is correct in spirit. The smooth activation is an additional
implementation choice that is not stated in the paper excerpt.

\section*{Closest-Point Detection}

\subsection*{Paper}

Paper Section 4.4 says that the original closest-point detection chooses the
smallest sensed distance from current sensor readings.

\subsection*{Current code}

Current code does exactly that first:

\begin{itemize}[leftmargin=*]
  \item compute all sensed vectors: lines 294--301
  \item compute distances: line 302
  \item choose smallest distance: lines 303--305
\end{itemize}

\textbf{Status}: \textbf{Exact}

\subsection*{Code excerpt}

\begin{lstlisting}[language=Python]
if hasattr(obstacle, "sensed_vectors"):
    sensed_vectors = obstacle.sensed_vectors(
        state.position,
        self.config.sensing_samples_per_obstacle,
        self.config.sensing_angular_window,
    )
else:
    sensed_vectors = [obstacle.closest_vector(state.position)]
distances = [_norm(vec) for vec in sensed_vectors]
min_index = int(np.argmin(distances))
dist_to_obs = sensed_vectors[min_index]
min_distance = distances[min_index]
\end{lstlisting}

\section*{Closest-Point Averaging for Non-Unique Closest Points}

\subsection*{Paper}

Paper Section 4.4 states:

\begin{itemize}[leftmargin=*]
  \item select several sensed distances smaller than threshold \(\delta r\),
  \item average them,
  \item for concave/non-unique cases, use the average if it is smaller in
  magnitude than the single closest sensed position,
  \item for convex cases, revert to the standard closest point.
\end{itemize}

\subsection*{Current code}

The current code follows exactly that logic:

\begin{itemize}[leftmargin=*]
  \item collect all vectors within the \(\delta r\) band: lines 307--309
  \item average them: line 310
  \item if the average is shorter, use it: lines 315--318
  \item otherwise use the single closest vector: lines 317--318
\end{itemize}

\textbf{Status}: \textbf{Exact in logic}

\subsection*{Code excerpt}

\begin{lstlisting}[language=Python]
close_vectors = [
    vec for vec, distance in zip(sensed_vectors, distances) if distance <= min_distance + self.config.delta_r
]
averaged_vector = np.mean(np.array(close_vectors, dtype=float), axis=0)

if _norm(averaged_vector) < min_distance:
    used_vector = averaged_vector
else:
    used_vector = dist_to_obs
\end{lstlisting}

\section*{Goal Relaxation}

\subsection*{Paper}

Paper Section 4.4 says goal relaxation is introduced for large obstacles where
following the boundary may increase the distance to the goal beyond the initial
goal distance. The paper description is qualitative in the Frontiers article:

\begin{itemize}[leftmargin=*]
  \item decrease \(F_g\) when the robot is close to the obstacle and the goal is
  still occluded,
  \item increase \(F_g\) when the obstacle no longer blocks the goal.
\end{itemize}

The paper explicitly points back to earlier work for details.

\subsection*{Current code}

The current implementation contains a legacy goal-relaxation weight:

\begin{equation}
k_{\text{GR}} =
\left(1-e^{-d_o/1.5}\right)
\left(1-\cos\theta\right)
\exp\left(-\frac{d_g-d_{g,0}}{0.1}\right)
\exp\left(-\frac{d_g-d_{g,\min}}{0.1}\right)
\end{equation}

This logic appears in lines 231--256 and is only active when
\texttt{use\_legacy\_goal\_relaxation = True} (line 269).

\subsection*{Code excerpt}

\begin{lstlisting}[language=Python]
goal_mag = _norm(goal_vector)
obs_mag = _norm(obs_vector)
if goal_mag < EPS or obs_mag >= self.config.r_l or obs_mag < EPS:
    self.min_distance_to_goal = min(self.min_distance_to_goal, goal_mag)
    return 1.0

if self.initial_goal_distance is None:
    self.initial_goal_distance = goal_mag
self.min_distance_to_goal = min(self.min_distance_to_goal, goal_mag)

cos_angle = float(np.dot(goal_vector, obs_vector) / max(goal_mag * obs_mag, EPS))
cos_angle = max(-1.0, min(1.0, cos_angle))
weight = 1.0 - cos_angle

if self.initial_goal_distance is None or goal_mag <= self.initial_goal_distance:
    new_weight = 1.0
else:
    new_weight = math.exp(-(goal_mag - self.initial_goal_distance) / 0.1)

weight_exit = math.exp(-(goal_mag - self.min_distance_to_goal) / 0.1)
distance_scale = 1.0 - math.exp(-obs_mag / 1.5)
return distance_scale * weight * new_weight * weight_exit
\end{lstlisting}

\subsection*{Comparison}

\textbf{Status}: \textbf{Approximate / inherited from earlier ROS work}

This is not a direct equation-level implementation from the 2022 Frontiers
paper. It is a pragmatic reuse of an earlier goal-relaxation mechanism from the
legacy codebase.

\section*{Local Sensing Assumption}

\subsection*{Paper}

The paper repeatedly states that only local sensory information is assumed, and
that no prior map is required.

\subsection*{Current code}

The current code honors that spirit through local obstacle samples, but the
obstacles themselves are analytic objects:

\begin{itemize}[leftmargin=*]
  \item circles sample local boundary points: lines 52--63
  \item polygons sample edge-local points: lines 98--114
\end{itemize}

\textbf{Status}: \textbf{Approximate}

This is a clean local-sensing simulation abstraction, not the same as the ROS
range-sensor or point-cloud pipeline used in the paper experiments.

\section*{Dynamics and Actuation Assumptions}

\subsection*{Paper}

The paper explicitly assumes non-saturatable actuators.

\subsection*{Current code}

The current simulator clips both acceleration and speed:

\begin{itemize}[leftmargin=*]
  \item acceleration clipping: lines 581--582
  \item speed clipping: lines 599--600
\end{itemize}

\textbf{Status}: \textbf{Different assumption}

This makes the simulation safer and numerically easier to manage, but it is not
the same assumption as the paper.

\section*{Point-Robot Parameter Table}

Paper Table 1 gives the following point-robot parameters:

\begin{itemize}[leftmargin=*]
  \item \(c = 10\)
  \item \(c_{\perp} = 20\)
  \item \(K_P = 0.04\)
  \item \(K_D = 0.5\)
  \item \(r_l = 3\)
  \item \(r_{la} = 2\)
  \item \(\epsilon = 3 \times 10^{-6}\)
\end{itemize}

The code path intended to mirror these values is:

\begin{itemize}[leftmargin=*]
  \item \texttt{make\_paper\_faithful\_config()} at lines 225--243
\end{itemize}

\textbf{Status}: \textbf{Exact for the numeric table, not yet exact for the full behaviour}

\section*{Overall Assessment}

\begin{longtable}{|p{0.34\linewidth}|p{0.16\linewidth}|p{0.34\linewidth}|}
\hline
\textbf{Component} & \textbf{Status} & \textbf{Comment} \\
\hline
High-level split \(u = F_g + F_b + F_a\) & Exact & Architecture matches the paper. \\
\hline
Closest-point detection & Exact & Single minimum-distance reading is implemented directly. \\
\hline
Section 4.4 closest-point averaging & Exact in logic & One of the strongest matches in the current code. \\
\hline
PD goal attraction & Approximate & Present, but wrapped in configurable modes and optional GR scaling. \\
\hline
Geometric control path & Approximate & Implemented in spirit, not yet verified from the cited controller derivation. \\
\hline
Boundary-following magnetic field \(F_b\) & Approximate & Current implementation captures intent, but the exact field law is still not reconstructed. \\
\hline
Collision-avoidance field \(F_a\) & Approximate & Same situation as \(F_b\). \\
\hline
Goal relaxation & Approximate / inherited & Current GR is inherited from earlier code, not reconstructed from the 2022 paper itself. \\
\hline
Local sensing assumption & Approximate & Simulation abstraction is local, but not the same sensing pipeline as the ROS work. \\
\hline
Actuation assumptions & Different & Paper assumes non-saturatable actuators; code clips acceleration and speed. \\
\hline
\end{longtable}

\section*{Practical Conclusion}

The current implementation is closest to the paper in:

\begin{itemize}[leftmargin=*]
  \item overall architecture,
  \item thresholded obstacle-field activation,
  \item closest-point and closest-point-averaging logic.
\end{itemize}

The current implementation is farthest from the paper in:

\begin{itemize}[leftmargin=*]
  \item the exact magnetic-field equations for \(F_b\) and \(F_a\),
  \item the exact goal-control choice used in each paper experiment,
  \item the exact sensing and actuation assumptions used in the paper setup.
\end{itemize}

Therefore, the current repository should still be described as a
\textbf{paper-inspired clean reimplementation}, not yet a verified reproduction
of the 2022 point-robot algorithm.

\end{document}
